% GENERATED FILE. Edit canonical lessons, then run npm run generate:books.
% canonical-source-hash: fnv1a64:2b95d9ae410944fe
% canonical-lessons: ZH-C505-wudao, ZH-C505-wuhui, ZH-W505-wu-dance, ZH-C505-wuban, ZH-C505-wutai, ZH-C505-pailian

\chapter{Dancing}
\label{ch:zh-dancing}
\hlchaptermodality{505}

\begin{chapteropening}
\textbf{By the end of this chapter:} \emph{I can talk about dance, a dance party, a dance partner, the stage and rehearsals.}

Complete the last lesson of chapter 505: I can talk about dance, a dance party, a dance partner, the stage and rehearsals.
\end{chapteropening}

\section[wǔdǎo]{wǔdǎo — the art of dance}

\label{lesson:ZH-C505-wudao}



\begin{quote}
Before the new one: say the Mandarin for a folk song, then the Mandarin for to perform a song (on stage).
\end{quote}

\subsection*{You'll want to know: wǔdǎo}
\emph{wǔdǎo} — \textquotedblleft{}the art of dance\textquotedblright{}.

\begin{quote}
\emph{wǒ nǚ'ér xué le wǔ nián wǔdǎo} — My daughter has studied dance for five years.
\end{quote}

Its characters come later; for now it is a word for the ear.

Two third tones meet in this word: wǔ is said with a rising tone before dǎo, and the spelling keeps both marks.

\begin{cousinweb}
\emph{wǔdǎo} is built from \emph{wǔ} 'to dance' and \emph{dǎo} 'to tread', both inherited from Middle Chinese, \emph{mjuX} and \emph{dawH} (Guangyun, in Baxter's transcription), which glosses the first 'singing and dancing' and the second 'to tread'. The modern third tone of dǎo is irregular: \emph{dawH} would regularly give dào.
\end{cousinweb}

\begin{grammarlens}[title={What you've built}]
One more word for saying how it was.
\end{grammarlens}

\subsection*{Your turn}
\begin{itemize}
\raggedright
  \item \emph{Say it:} \emph{wǔdǎo}
  \item \emph{Say it:} \emph{wǔdǎo}, once more
  \item \emph{Say it:} \emph{wǒ nǚ'ér xué le wǔ nián wǔdǎo} — My daughter has studied dance for five years.
  \item \emph{Recall:} say the Mandarin for a folk song, then the Mandarin for to perform a song (on stage), then say \emph{wǔdǎo} again
\end{itemize}

\begin{tcolorbox}[breakable,colback=teal!4,colframe=teal!35!black,title={Before you move on}]
What does wǔdǎo mean? (\textquotedblleft{}The art of dance\textquotedblright{}.) Say it once more.
\end{tcolorbox}
\section[\texorpdfstring{\zh{舞会} (wǔhuì)}{舞会 (wǔhuì)}]{\zh{舞会} (wǔhuì) — dance party; ball}

\label{lesson:ZH-C505-wuhui}



\begin{quote}
Before the new one: say the Mandarin for to perform a song (on stage), then the Mandarin for the art of dance.
\end{quote}

\subsection*{You'll want to know: \zh{舞会}}
\textbf{\zh{舞会}} — \emph{wǔhuì} — \textquotedblleft{}dance party; ball\textquotedblright{}.

\begin{quote}
\emph{zhè shì wǒ dìyī cì cānjiā wǔhuì} — This is the first time I've been to a dance.
\end{quote}

The new character, in pieces: \zh{舞} — fourteen strokes: on top eight, a short stroke falling left, two bars with four short verticals through the lower one, and a long bar across; below, on the left the top of \zh{多} \emph{duō}, three strokes, and on the right three strokes. The next lesson puts it in your hand.

\begin{cousinweb}
\textbf{\zh{舞会}} (\emph{wǔhuì}) is built from \zh{舞} 'to dance' and \zh{会} 'a gathering', both inherited from Middle Chinese, \emph{mjuX} and \emph{hwajH} (Guangyun, in Baxter's transcription), which glosses the first 'singing and dancing' and the second 'to meet'.
\end{cousinweb}

\begin{grammarlens}[title={What you've built}]
One more word for saying how it was.
\end{grammarlens}

\subsection*{Your turn}
\begin{itemize}
\raggedright
  \item \emph{Say it:} \emph{wǔhuì}
  \item \emph{Say it:} \emph{wǔhuì}, once more
  \item \emph{Say it:} \emph{zhè shì wǒ dìyī cì cānjiā wǔhuì} — This is the first time I've been to a dance.
  \item \emph{Recall:} say the Mandarin for to perform a song (on stage), then the Mandarin for the art of dance, then say \emph{wǔhuì} again
\end{itemize}

\begin{tcolorbox}[breakable,colback=teal!4,colframe=teal!35!black,title={Before you move on}]
What does \zh{舞会} mean? (\textquotedblleft{}Dance party; ball\textquotedblright{}.) Say it once more.
\end{tcolorbox}
\section[\texorpdfstring{\zh{舞} (wǔ)}{舞 (wǔ)}]{\zh{舞} — eight strokes on top, then two pieces below}

\label{lesson:ZH-W505-wu-dance}



\begin{quote}
The word first, then the character inside it.

\begin{itemize}
\raggedright
  \item \emph{Recall:} say \emph{wǔhuì} — \textbf{R1}, one lesson back, before it decays
\end{itemize}

\begin{quote}
\zh{多} \zh{会}
\end{quote}

\emph{Read it:} each one aloud, with nothing said first — \emph{duō}, \emph{huì}
\end{quote}

\subsection*{Script — eight strokes on top, then two pieces below}
\hlblockfigure{\detokenize{figures/ZH-W505-wu-dance-filmstrip.pdf}}{How \zh{舞} is written, stroke by stroke}

\begin{quote}
\zh{舞}
\end{quote}

\textbf{The top, eight strokes.} A short stroke falling to the left. A short bar, then a longer bar under it. Four short verticals through it, from left to right. Then a long bar across the whole width.

\textbf{The lower left, three strokes}, as on top of \zh{多}: a stroke falling to the left, a bar that turns and falls to the left, and a dot.

\textbf{The lower right, three strokes}: a bar, a stroke that drops and turns to the right, then a long vertical down through them.

8 + 3 + 3 = fourteen.

\subsection*{Your turn}
\begin{itemize}
\raggedright
  \item \emph{Write it:} \zh{舞}, all fourteen, with the model in front of you
  \item \emph{Cover:} write it again from memory, and count the pieces, not the strokes
\end{itemize}

\begin{tcolorbox}[breakable,colback=teal!4,colframe=teal!35!black,title={Before you move on}]
How many strokes does \zh{舞} have? (\textbf{Fourteen}.) How many short verticals in the top? (\textbf{Four}.) Which character you know has the lower left piece on top? (\textbf{\zh{多}}.)

Source: \href{https://raw.githubusercontent.com/chanind/hanzi-writer-data/68d10a4b21150cae5e1ebbd223eed289cf32d90c/data/\%E8\%88\%9E.json}{Hanzi Writer Data \zh{舞}, pinned stroke paths and medians}.
\end{tcolorbox}
\section[wǔbàn]{wǔbàn — dance partner}

\label{lesson:ZH-C505-wuban}



\begin{quote}
Before the new one: say the Mandarin for the art of dance, then the Mandarin for a dance party.
\end{quote}

\subsection*{You'll want to know: wǔbàn}
\emph{wǔbàn} — \textquotedblleft{}dance partner\textquotedblright{}.

\begin{quote}
\emph{tā shì wǒ zuì hǎo de wǔbàn} — She's my best dance partner.
\end{quote}

Its characters come later; for now it is a word for the ear.

\begin{cousinweb}
\emph{wǔbàn} is built from \zh{舞} 'to dance' and \emph{bàn} 'a companion', both inherited from Middle Chinese, \emph{mjuX} and \emph{banX} (Guangyun, in Baxter's transcription), which glosses the second 'a companion'.
\end{cousinweb}

\begin{grammarlens}[title={What you've built}]
One more word for saying how it was.
\end{grammarlens}

\subsection*{Your turn}
\begin{itemize}
\raggedright
  \item \emph{Say it:} \emph{wǔbàn}
  \item \emph{Say it:} \emph{wǔbàn}, once more
  \item \emph{Say it:} \emph{tā shì wǒ zuì hǎo de wǔbàn} — She's my best dance partner.
  \item \emph{Recall:} say the Mandarin for the art of dance, then the Mandarin for a dance party, then say \emph{wǔbàn} again
\end{itemize}

\begin{tcolorbox}[breakable,colback=teal!4,colframe=teal!35!black,title={Before you move on}]
What does wǔbàn mean? (\textquotedblleft{}Dance partner\textquotedblright{}.) Say it once more.
\end{tcolorbox}
\section[wǔtái]{wǔtái — stage (in a theatre)}

\label{lesson:ZH-C505-wutai}



\begin{quote}
Before the new one: say the Mandarin for a dance party, then the Mandarin for a dance partner.
\end{quote}

\subsection*{You'll want to know: wǔtái}
\emph{wǔtái} — \textquotedblleft{}stage (in a theatre)\textquotedblright{}.

\begin{quote}
\emph{nàge wǔtái bǐ wǒmen xuéxiào de dà} — That stage is bigger than our school's.
\end{quote}

Its characters come later; for now it is a word for the ear.

\begin{cousinweb}
\emph{wǔtái} is built from \zh{舞} 'to dance' and \emph{tái} 'a platform', both inherited from Middle Chinese, \emph{mjuX} and \emph{doj} (Guangyun, in Baxter's transcription), which glosses the second as a raised square of earth.
\end{cousinweb}

\begin{grammarlens}[title={What you've built}]
One more word for saying how it was.
\end{grammarlens}

\subsection*{Your turn}
\begin{itemize}
\raggedright
  \item \emph{Say it:} \emph{wǔtái}
  \item \emph{Say it:} \emph{wǔtái}, once more
  \item \emph{Say it:} \emph{nàge wǔtái bǐ wǒmen xuéxiào de dà} — That stage is bigger than our school's.
  \item \emph{Recall:} say the Mandarin for a dance party, then the Mandarin for a dance partner, then say \emph{wǔtái} again
\end{itemize}

\begin{tcolorbox}[breakable,colback=teal!4,colframe=teal!35!black,title={Before you move on}]
What does wǔtái mean? (\textquotedblleft{}Stage (in a theatre)\textquotedblright{}.) Say it once more.
\end{tcolorbox}
\section[páiliàn]{páiliàn — to rehearse}

\label{lesson:ZH-C505-pailian}



\begin{quote}
Before the new one: say the Mandarin for a dance partner, then the Mandarin for a stage (in a theatre).
\end{quote}

\subsection*{You'll want to know: páiliàn}
\emph{páiliàn} — \textquotedblleft{}to rehearse\textquotedblright{}.

\begin{quote}
\emph{wǒmen páiliàn le yí ge yuè} — We rehearsed for a month.
\end{quote}

Its characters come later; for now it is a word for the ear.

\begin{cousinweb}
\emph{páiliàn} is built from \emph{pái} 'to line up' and \emph{liàn} 'to practise', both inherited from Middle Chinese, \emph{beaj} and \emph{lenH} (Guangyun, in Baxter's transcription), which glosses the first 'to push' and the second 'white silk'; the senses 'to line up' and 'to practise' are not given there.
\end{cousinweb}

\begin{grammarlens}[title={What you've built}]
One more word for saying how it was.
\end{grammarlens}

\subsection*{Your turn}
\begin{itemize}
\raggedright
  \item \emph{Say it:} \emph{páiliàn}
  \item \emph{Say it:} \emph{páiliàn}, once more
  \item \emph{Say it:} \emph{wǒmen páiliàn le yí ge yuè} — We rehearsed for a month.
  \item \emph{Recall:} say the Mandarin for a dance partner, then the Mandarin for a stage (in a theatre), then say \emph{páiliàn} again
\end{itemize}

\begin{tcolorbox}[breakable,colback=teal!4,colframe=teal!35!black,title={Before you move on}]
What does páiliàn mean? (\textquotedblleft{}To rehearse\textquotedblright{}.) Say it once more.
\end{tcolorbox}
